\documentclass[11pt]{article}
\usepackage[a4paper,margin=25mm]{geometry}
\usepackage{amsmath,amssymb,amsthm,mathtools,booktabs}
\usepackage[T1]{fontenc}
\usepackage{lmodern,microtype,needspace}
\usepackage[hidelinks]{hyperref}
\hypersetup{pdftitle={Sharp boundary bounds and edit instability in uniquely dominated bipartite graphs},pdfauthor={}}
\newtheorem{theorem}{Theorem}[section]
\newtheorem{lemma}[theorem]{Lemma}
\newtheorem{proposition}[theorem]{Proposition}
\newtheorem{corollary}[theorem]{Corollary}
\theoremstyle{definition}
\newtheorem{definition}[theorem]{Definition}
\theoremstyle{remark}
\newtheorem{remark}[theorem]{Remark}
\newcommand{\epn}{\operatorname{epn}}
\newcommand{\e}{e}
\newcommand{\ceil}[1]{\left\lceil#1\right\rceil}
\newcommand{\floor}[1]{\left\lfloor#1\right\rfloor}
\newcommand{\setminusone}[2]{#1\setminus\{#2\}}
\setlength{\emergencystretch}{2em}
\clubpenalty=10000 \widowpenalty=10000
\title{Sharp boundary bounds and edit instability\\in uniquely dominated bipartite graphs}
\author{}
\date{5 October 2026\\\small Third version}
\begin{document}
\maketitle
\begin{abstract}
We determine the sharp edge bounds for bipartite graphs without isolated
vertices that have a unique minimum dominating set of size $\gamma\ge2$ and
order $3\gamma+1$ or $3\gamma+2$. The maxima are $\gamma(\gamma+7)/2$ and
$\lceil\gamma^2/2\rceil+5\gamma$, respectively, and all equality graphs are
classified. There is one equality class at the first order. At the second
order there are three classes when $\gamma=2$, one for odd $\gamma\ge3$,
and two for even $\gamma\ge4$.
We then study edge-edit distance, minimized over all vertex relabellings,
at order $3\gamma+1$. A graph with edge deficit $t$ is within $12\gamma t$
edits of the unique extremal graph. Connected examples require
$\Omega(\gamma t)$ edits for $1\le t\le\gamma/2-1$ when $\gamma$ is even.
In particular, relative edge deficit can tend to zero while normalized
edit distance stays bounded away from zero. The proofs are elementary,
using private-neighbour replacements, local edge deficits, and explicit
near-extremal constructions. Previously reported finite counterexamples
and construction subfamilies are recovered and explicitly credited.
\end{abstract}
\noindent\textbf{Keywords:} unique minimum dominating set; bipartite graph;
extremal size; private neighbour.
\par\smallskip
\noindent\textbf{2020 Mathematics Subject Classification:} 05C69, 05C35.

\section{Introduction}

All graphs in this paper are finite and simple. A set $D\subseteq V(G)$ is
\emph{dominating} if every vertex outside $D$ has a neighbour in $D$.
The domination number $\gamma(G)$ is the minimum cardinality of a dominating
set. A graph has a \emph{unique minimum dominating set} if exactly one dominating
set has cardinality $\gamma(G)$. Throughout, uniqueness refers to minimum
cardinality, not to inclusion-minimality. Write $n(G)=|V(G)|$ and
$\e(G)=|E(G)|$.

Fischermann, Rautenbach and Volkmann~\cite{FRV} investigated the maximum size
of graphs with a unique minimum dominating set. Fraboni and Shank~\cite{FS}
subsequently settled their conjectured size bound when the domination number
is two. Koch and Narayan~\cite{KN} considered the bipartite problem and proposed
a bound depending on the order and the domination number. They proved their
bound for domination number two and at order $n=3\gamma$.
Erlbacher's public artifacts~\cite{E} already give the $13$-vertex, $22$-edge
counterexample, an unbalanced family containing a $14$-vertex, $28$-edge
example, and finite extremal computations. The results below recover those
examples, determine the sharp maxima and equality classes at the next two
orders, and examine their quantitative edit stability.
Koch and Narayan also identify the unique extremal graph at
$(n,\gamma)=(10,3)$ in their Figure~1 discussion.

The order restriction is intrinsic. If $D$ is the unique minimum dominating
set of a graph with no isolated vertices, every vertex of $D$ has at least two
exterior private neighbours with respect to $D$. These neighbours are distinct
for distinct members of $D$, and hence $n(G)\ge3\gamma(G)$. We recall the short
argument in Lemma~\ref{lem:private}.

For integers $\gamma\ge2$ and $r\in\{1,2\}$, let $\mathcal B_{\gamma,r}$ be the class of bipartite
graphs with no isolated vertices, of order $3\gamma+r$, whose minimum
dominating set is unique and has cardinality $\gamma$. Write
$\mathcal B_\gamma=\mathcal B_{\gamma,1}$. No connectedness assumption is
imposed in these definitions.

\Needspace{10\baselineskip}
\begin{theorem}\label{thm:main}
For every integer $\gamma\ge2$,
\[
 \max_{G\in\mathcal B_\gamma}\e(G)=\frac{\gamma(\gamma+7)}2.
\]
The maximum is attained by a connected graph. For every $\gamma\ge2$, equality
holds if and only if
\[
 G\cong H\bigl(\ceil{\gamma/2},\floor{\gamma/2}\bigr),
\]
where $H(p,q)$ is the graph defined in Section~\ref{sec:construction}.
\end{theorem}

\begin{theorem}\label{thm:two}
For every integer $\gamma\ge2$,
\[
 \max_{G\in\mathcal B_{\gamma,2}}\e(G)
   =\ceil{\frac{\gamma^2}{2}}+5\gamma.
\]
The equality graphs, up to isomorphism, are exactly
\[
\begin{cases}
 H_2(1,1),\ F_0,\ F_1,&\gamma=2,\\
 H_2(k+1,k),&\gamma=2k+1\ge3,\\
 H_2(k,k),\ H_2(k+1,k-1),&\gamma=2k\ge4.
\end{cases}
\]
Here $H_2$ is defined in Section~\ref{sec:construction}, and the two
balanced eight-vertex exceptions $F_0,F_1$ in Section~\ref{sec:two-rigidity}.
Every equality graph is connected.
\end{theorem}

For graphs $G,H$ of the same order, define
\[
 d_{\mathrm{edit}}(G,H)
   =\min_{\phi:V(H)\to V(G)\ \mathrm{bijective}}
       |E(G)\mathbin{\triangle}\phi(E(H))|.
\]
Each edge addition or deletion costs one. The bijection need not preserve
the displayed bipartition or a named dominating set.

\begin{theorem}\label{thm:stability}
Let $\gamma\ge2$ and $G\in\mathcal B_{\gamma,1}$, let
$H^*=H(\ceil{\gamma/2},\floor{\gamma/2})$, and put
$t=\gamma(\gamma+7)/2-\e(G)$. Then
\[
 d_{\mathrm{edit}}(G,H^*)
 \le \min\{\gamma(\gamma+7)-t,\ 12\gamma t\}.
\]
For every $k\ge2$ and $1\le s\le k-1$, there is a connected graph
$G(k,s)\in\mathcal B_{2k,1}$ with exact deficit $t=s$ such that
\[
 s(k+1)\le d_{\mathrm{edit}}(G(k,s),H(k,k))\le s(2k+1).
\]
\end{theorem}

Thus the dependence $\gamma t$ has the correct order for small absolute
edge deficit. It cannot generally be replaced by $\gamma\sqrt t+t$.
More strongly, the connected examples have relative edge deficit tending
to zero but normalized edit distance bounded away from zero; the precise
limits are given in Section~\ref{sec:obstruction}.


At order $3\gamma+1$, Conjecture~1 of~\cite{KN} specializes to
\begin{equation}\label{eq:sourcebound}
 \e(G)\le 2\gamma+2ab+2a+1,
 \qquad a=\ceil{\gamma/2},\quad b=\floor{\gamma/2}.
\end{equation}
The summation term in the conjecture is empty at this order. Since
\begin{equation}\label{eq:gap}
 \frac{\gamma(\gamma+7)}2-(2\gamma+2ab+2a+1)=b-1,
\end{equation}
Theorem~\ref{thm:main} has the following consequence.

\begin{corollary}\label{cor:counter}
For every $\gamma\ge4$, there is a connected bipartite graph on $3\gamma+1$
vertices with a unique minimum dominating set of cardinality $\gamma$ that
violates~\eqref{eq:sourcebound}. The excess is $\floor{\gamma/2}-1$.
In particular, at $\gamma=4$ the extremal graph has $13$ vertices and $22$
edges, whereas~\eqref{eq:sourcebound} gives $21$.
\end{corollary}

The additional vertex explains the change from the threshold case $n=3\gamma$.
It may be adjacent to several members of the unique minimum dominating set.
Our extremal graph uses this possibility. The proof of the upper bound keeps
track of whether the additional vertex has one or several such neighbours;
the equality argument then forces the whole construction.

\section{The extremal constructions}\label{sec:construction}

\begin{definition}\label{def:H}
For integers $p,q\ge1$ and $r\in\{1,2\}$, let $H_r(p,q)$ have pairwise
 distinct vertices
\[
 X=\{x_i:1\le i\le p\},\qquad Y=\{y_j:1\le j\le q\},
\]
\[
 U_i=\{u_i^0,u_i^1\},\qquad V_j=\{v_j^0,v_j^1\},\qquad
 Z=\{z_1,\ldots,z_r\}.
\]
Its edges are exactly $x_i u_i^\varepsilon$, $y_j v_j^\varepsilon$,
$u_i^\varepsilon v_j^0$, $z_tu_i^\varepsilon$ and $z_ty_j$, for every
admissible index and $\varepsilon\in\{0,1\}$. Write $H(p,q)=H_1(p,q)$,
and write $z$ for its single member of $Z$.
\end{definition}

The graph is bipartite, with parts
\begin{equation}\label{eq:parts}
 X\cup\bigcup_jV_j\cup Z
 \qquad\text{and}\qquad
 Y\cup\bigcup_iU_i.
\end{equation}
It is connected: each $z_t$ meets all of $Y\cup\bigcup_iU_i$, each $x_i$
meets $U_i$, and each member of $V_j$ meets $y_j$.

\begin{proposition}\label{prop:H}
The graph $H_r(p,q)$ has order $3(p+q)+r$, domination number $p+q$, and
unique minimum dominating set $X\cup Y$. Its size is
\begin{equation}\label{eq:Hcount}
 \e(H_r(p,q))=2(p+q)+2pq+r(2p+q).
\end{equation}
\end{proposition}

\begin{proof}
The vertex and edge counts follow from the definition. Clearly
$D=X\cup Y$ dominates. Every dominating set $K$ meets each closed
neighbourhood $N[x_i]=\{x_i,u_i^0,u_i^1\}$ and each
$N[v_j^1]=\{y_j,v_j^1\}$. These $p+q$ sets are pairwise disjoint, so
$|K|\ge p+q$.

If $|K|=p+q$, it chooses exactly one member of each of these sets and
nothing else. Thus every member of $Z$, and every $v_j^0$, is absent.
Choosing $u_i^\varepsilon$ instead of $x_i$ leaves its mate undominated:
the mate's neighbours are $x_i$, the members of $Z$, and the vertices
$v_j^0$. Hence all $x_i$ are selected and all $U_i$ vertices are absent.
Choosing $v_j^1$ instead of $y_j$ now leaves $v_j^0$ undominated, forcing
all $y_j$. Thus $K=D$.
\end{proof}

For $a=\ceil{\gamma/2}$ and $b=\floor{\gamma/2}$, direct calculation gives
\begin{align}
 \e(H(a,b))&=2ab+4a+3b=\frac{\gamma(\gamma+7)}2,\label{eq:attain}\\
 \e(H_2(a,b))&=2ab+6a+4b
     =\ceil{\frac{\gamma^2}{2}}+5\gamma.\label{eq:attain2}
\end{align}
These graphs supply the connected lower bounds for both theorems.

The subfamily $H_2(p,1)$ is the two-extra-vertex case of Erlbacher's
unbalanced-split construction~\cite{E}, after interchanging the names of
the center sides. In particular, $H_2(3,1)$ is his $14$-vertex example.
For $q>1$, each member of $Z$ meets several centers; this shared-center
feature already occurs in his $13$-vertex counterexample, which is $H(2,2)$.
We use these constructions to establish sharpness, without claiming
priority for those previously reported examples or mechanisms.

\section{Private neighbours and local replacements}

For a dominating set $D$ and $d\in D$, define
\[
 \epn(d,D)=\{u\in V(G)\setminus D:N(u)\cap D=\{d\}\}.
\]
We use closed neighbourhoods only when square brackets are displayed.

\Needspace{6\baselineskip}
\begin{lemma}\label{lem:private}
If $G$ has no isolated vertices and $D$ is its unique minimum dominating set,
then $|\epn(d,D)|\ge2$ for every $d\in D$.
\end{lemma}

\begin{proof}
If $\epn(d,D)=\{u\}$, the set $(D\setminus\{d\})\cup\{u\}$ dominates $G$:
$u$ dominates itself and $d$, while every other vertex outside $D$ retains a
neighbour in $D\setminus\{d\}$. This contradicts uniqueness.

If $\epn(d,D)=\varnothing$ and $d$ has a neighbour in $D$, then
$D\setminus\{d\}$ is a smaller dominating set. Otherwise $d$ has no neighbour
in $D$. The no-isolate assumption gives a neighbour $u$ outside $D$, and
$(D\setminus\{d\})\cup\{u\}$ is again a different dominating set of the same
cardinality. All cases are impossible.
\end{proof}

The local lemmas below include every possible edge between the two centers.
Their use is governed by the following replacement principle. If an induced
subgraph $F$ contains $x,y\in D$ and all other members of $D$ lie outside $F$,
and $D\setminus\{x,y\}$ dominates every vertex outside $F$, then any
dominating set $S$ of $F$ of size at most two extends to
$(D\setminus\{x,y\})\cup S$. Thus, whenever $\{x,y\}$ dominates $F$, it
is its unique dominating set of size at most two. In each application we
verify the outside-domination condition explicitly; cross-cell edges cannot
obstruct the resulting domination.
The six-vertex argument occurs in the proofs of Theorem~1 of~\cite{FRV} and
Theorem~12 of~\cite{KN}; we recall its equality patterns.

\begin{lemma}\label{lem:six}
Let $F$ be a bipartite graph with parts $\{x\}\cup V$ and $\{y\}\cup U$,
where $|U|=|V|=2$. Assume that $x$ is adjacent to both vertices of $U$ and
$y$ to both vertices of $V$. If $\{x,y\}$ is the unique dominating set of
cardinality at most two, then at most two edges occur among $xy$ and the
$U$--$V$ edges. If exactly two occur, either $xy$ is absent and the
$U$--$V$ edges form a two-edge star, or $xy$ is present and there is exactly
one $U$--$V$ edge.
\end{lemma}

\begin{proof}
Two disjoint $U$--$V$ edges give an alternative dominating pair: select
opposite endpoints, one from each edge. If $xy$ is absent, three
$U$--$V$ edges contain two disjoint edges. Thus there are at most two, and a
two-edge configuration must be a star.

Suppose $xy$ is present. Two $U$--$V$ edges either are disjoint or form a
star. A star centered at $u\in U$ gives the dominating pair $\{x,u\}$;
a star centered at $v\in V$ gives $\{y,v\}$. Consequently at most one
$U$--$V$ edge can accompany $xy$.
\end{proof}

\begin{lemma}\label{lem:seven}
Let $F$ have the same form as in Lemma~\ref{lem:six}, with $|U|=2$ and
$|V|=s\in\{3,4\}$. If $\{x,y\}$ is its unique dominating set of size at
most two, then at most $2s-2$ optional edges occur among $xy$ and $U$--$V$.
\end{lemma}

\begin{proof}
If $xy$ is absent and there are at least $2s-1$ cross-edges, one $U$ row
is full and the other has a neighbour $v$. The full row vertex $u$ and $v$
form an alternative dominating pair. If $xy$ is present, a $V$ vertex
adjacent to both members of $U$ gives the alternative $\{y,v\}$.
Hence every column has at most one cross-edge, leaving at most
$s+1\le2s-2$ optional edges in total.
\end{proof}

\begin{lemma}\label{lem:seven-equality}
In Lemma~\ref{lem:seven} with $s=3$, equality holds exactly when $xy$ is
absent and the cross-edges form a $K_{2,2}$ on $U$ and two members of $V$,
leaving the third column empty.
\end{lemma}

\begin{proof}
If $xy$ is present at equality, there are three cross-edges and no full
column. Each column therefore contains one edge. If one row contains all
three, $\{x,u\}$ dominates. Otherwise the row degrees are two and one;
the degree-two row vertex together with the other row's neighbour dominates.
Both contradict uniqueness.

With $xy$ absent, four cross-edges have row degrees $(3,1)$ or $(2,2)$.
A full row vertex and a neighbour of the other row rule out $(3,1)$.
For $(2,2)$ with different row neighbourhoods, select $u_0$ and the unique
member of $N(u_1)\cap V$ outside $N(u_0)$. They dominate the whole cell.
Thus both rows meet the same two columns.

Conversely, let $c$ be the empty column. A dominating pair must meet both
disjoint sets $N[x]=\{x,u_0,u_1\}$ and $N[c]=\{c,y\}$, and selects nothing
else. Choosing a $U$ vertex leaves its mate undominated, forcing $x$;
then choosing $c$ leaves the active columns undominated, forcing $y$.
No singleton can meet both disjoint sets.
\end{proof}

\begin{lemma}\label{lem:eight33}
Let $F$ have the same form, now with $|U|=|V|=3$. If $\{x,y\}$ is its unique
dominating set of size at most two, then at most six optional edges occur.
\end{lemma}

\begin{proof}
Suppose at least seven optional edges occur. If $xy$ is absent, the $3$ by
$3$ cross-edge matrix has at most two missing entries. A full row and a
full column therefore exist; their vertices form an alternative dominating
pair. If $xy$ is present, there are at least six cross-edges. A full row
gives $\{x,u\}$, and a full column gives $\{y,v\}$. Otherwise every row and
column has degree at most two. Six edges force every degree to equal two,
so the missing entries form a perfect matching. The endpoints of any
missing edge dominate both centers and every private vertex, again a
contradiction.
\end{proof}

These elementary arguments give the optional-edge bounds $2,4,6,6$ for
private-set sizes $(2,2),(2,3),(2,4),(3,3)$ respectively. The exact checks
in Appendix~\ref{app:checks} are supplementary.

\section{One residual vertex}\label{sec:upper}

Let $G\in\mathcal B_\gamma$, and let $D$ be its unique minimum dominating set.
By Lemma~\ref{lem:private}, we may select two exterior private neighbours for
each vertex of $D$. The selected pairs are disjoint and account for $2\gamma$
vertices. Since $n(G)=3\gamma+1$, exactly one further vertex $z$ remains.

Choose the names of the bipartition classes so that $z\in A$, and write
\[
 D\cap A=\{x_1,\ldots,x_p\},\qquad
 D\cap B=\{y_1,\ldots,y_q\},\qquad p+q=\gamma.
\]
Let $U_i\subseteq B$ be the selected private pair of $x_i$, and
$V_j\subseteq A$ the selected private pair of $y_j$. Set
\[
 J=N(z)\cap D\subseteq\{y_1,\ldots,y_q\},\qquad \ell=|J|.
\]
Because $D$ dominates $z$, we have $\ell\ge1$, and in particular $q\ge1$.

Privacy and bipartiteness give a complete list of possible edges:
\begin{enumerate}
\item the $2\gamma$ selected center-to-private edges;
\item $x_iy_j$ and the edges between $U_i$ and $V_j$, for each pair $(i,j)$;
\item the edges from $z$ to $J$ and to $\bigcup_iU_i$.
\end{enumerate}
In particular, edges inside $D$ have not been discarded: they are precisely
the possible edges $x_iy_j$ in the second class.

\subsection{The additional vertex has several dominators}

Suppose $\ell\ge2$. For each $i,j$, consider the induced six-vertex graph
\[
 G[K_{ij}],\qquad K_{ij}=\{x_i,y_j\}\cup U_i\cup V_j.
\]
The set $D\setminus\{x_i,y_j\}$ dominates every vertex outside $K_{ij}$.
Each selected private vertex outside the cell retains its own center;
the vertex $z$ retains a neighbour in $J$ because only one $y$-center has
been removed.

If $G[K_{ij}]$ had a dominating set $S$ of cardinality at most two other
than $\{x_i,y_j\}$, then
\[
 (D\setminus\{x_i,y_j\})\cup S
\]
would be a dominating set of $G$ of cardinality at most $\gamma$, distinct
from $D$ when its cardinality is $\gamma$. This contradicts minimum
cardinality or uniqueness. Thus Lemma~\ref{lem:six} applies to every cell.
The optional edge sets of these cells are disjoint, and
$\deg(z)\le2p+q$. Consequently
\begin{equation}\label{eq:multiple}
 \e(G)\le 2\gamma+2pq+2p+q.
\end{equation}
When $p=0$, the same bound follows directly from the edge list, with an
empty sum over cells.

\subsection{The additional vertex has one dominator}

Suppose $J=\{y_{j_0}\}$. For every $j\ne j_0$, the preceding six-vertex
replacement still applies, since $y_{j_0}$ is retained and dominates $z$.
For each $i$ consider instead
\[
 L_i=\{x_i,y_{j_0}\}\cup U_i\cup V_{j_0}\cup\{z\}.
\]
In $G[L_i]$, the center $y_{j_0}$ has the three private vertices
$V_{j_0}\cup\{z\}$. Every vertex outside $L_i$ is dominated by
$D\setminus\{x_i,y_{j_0}\}$, so the same replacement argument proves that
Lemma~\ref{lem:seven} applies.

The optional edges of $G[L_i]$ are $x_i y_{j_0}$ and the edges from $U_i$
to $V_{j_0}\cup\{z\}$. These edge sets are disjoint for different $i$.
The edge $zy_{j_0}$ is counted once separately. Therefore
\begin{align}
 \e(G)&\le 2\gamma+1+2p(q-1)+4p\notag\\
       &=2\gamma+1+2pq+2p\notag\\
       &\le2\gamma+2pq+2p+q.\label{eq:single}
\end{align}
If $p=0$, only the $2\gamma$ fixed private edges and $zy_{j_0}$ occur,
so the displayed inequalities remain valid.

\subsection{Optimizing the bipartition}

Put $d=p-q$. Since $p+q=\gamma$, both cases yield
\begin{equation}\label{eq:opt}
 \e(G)\le 2\gamma+2pq+2p+q
 =\frac{\gamma(\gamma+7)}2-\frac{d(d-1)}2
 \le\frac{\gamma(\gamma+7)}2.
\end{equation}
The final inequality holds because $d$ is an integer. Equality in the
numerical optimization requires $d\in\{0,1\}$. The parity of $\gamma$
then forces
\begin{equation}\label{eq:balanced}
 p=\ceil{\gamma/2},\qquad q=\floor{\gamma/2}.
\end{equation}
Together with~\eqref{eq:attain}, this proves the extremal-value assertion of
Theorem~\ref{thm:main} for every $\gamma\ge2$.

\section{Equality with one residual vertex}\label{sec:rigidity}

\subsection{Domination number at least four}

Suppose now that $\gamma\ge4$ and
$\e(G)=\gamma(\gamma+7)/2$. Equation~\eqref{eq:balanced} gives
$q=\floor{\gamma/2}\ge2$. The final inequality in~\eqref{eq:single}
loses $q-1$ edges, so the case $\ell=1$ cannot give equality.
Every inequality in~\eqref{eq:multiple} must therefore be an equality.
It follows that $z$ is adjacent to every $y_j$ and to every vertex of
$\bigcup_iU_i$, and each six-vertex cell has exactly two optional edges.

We first show that $D$ is independent. If $x_iy_j$ were an edge, then
\[
 (D\setminus\{x_i\})\cup\{z\}
\]
would be a different dominating set of cardinality $\gamma$.
Indeed, $x_i$ is dominated by the retained $y_j$, the pair $U_i$ is
dominated by $z$, and every other private pair retains its center.
Thus no edge $x_iy_j$ occurs.

Lemma~\ref{lem:six} now shows that the two edges between $U_i$ and $V_j$
form a star. This star cannot be centered at $u\in U_i$, because
\[
 (D\setminus\{x_i,y_j\})\cup\{u,z\}
\]
would dominate $G$. The vertex $u$ dominates $x_i$ and both members of
$V_j$, while $z$ dominates $y_j$ and both members of $U_i$.
All other private pairs retain their centers. Therefore both members of
$U_i$ are adjacent to exactly one vertex of $V_j$; denote it by $v_j^{(i)}$.

Fix $j$. Suppose that the selected vertex $v_j^{(i)}$ is not constant as
$i$ varies. Choose one vertex $u_i\in U_i$ for every $i$, and define
\[
 K=\{z\}\cup\{u_i:1\le i\le p\}
       \cup\bigl((D\cap B)\setminus\{y_j\}\bigr).
\]
Then $|K|=1+p+(q-1)=\gamma$. Every $x_i$ is dominated by $u_i$;
every vertex in $\bigcup_iU_i$ and every $y$-center is dominated by $z$.
The two vertices of $V_j$ are dominated because both possible choices of
$v_j^{(i)}$ occur. Every other pair $V_\ell$ retains $y_\ell$.
Hence $K$ is a second minimum dominating set, a contradiction.

For each $j$, rename the members of $V_j$ so that their common selected
vertex is $v_j^0$. The complete edge list is now exactly that of $H(p,q)$.
Equation~\eqref{eq:balanced} proves the equality assertion in
Theorem~\ref{thm:main} for $\gamma\ge4$.

\subsection{Domination numbers two and three}

At equality,~\eqref{eq:balanced} gives $q=1$ and $p=1$ or $2$.
Write $y$ for the sole $Y$-center and put $W=V_1\cup\{z\}$. All three
vertices of $W$ meet $y$. The graph has $2p+3$ fixed center/private edges,
and its optional edges partition into the $p$ blocks $x_i y$ and $U_i$--$W$.
Each seven-vertex cell is legitimate for replacement: every outside vertex
retains its own center, while $z$ is inside the cell. Equality forces four
optional edges in every block.

Lemma~\ref{lem:seven-equality} forbids all center edges and makes each $U_i$
complete to two of the three $W$ columns. For $p=1$ this is $H(1,1)$.
For $p=2$, suppose the omitted columns differ. The active column pairs
then cover $W$ and have one common member $w$. Choose $u_i\in U_i$.
The set $\{u_1,u_2,w\}$ dominates: the $u_i$ cover $x_i$ and all columns
between them, while $w$ covers both $U$ pairs and $y$. This is a different
dominating triple, a contradiction. The omitted columns therefore coincide,
and the graph is $H(2,1)$. This completes Theorem~\ref{thm:main} for all
$\gamma\ge2$. As noted above, the latter finite equality graph was already
identified in~\cite{KN}.

\section{Two residual vertices}\label{sec:two}

Let $G\in\mathcal B_{\gamma,2}$ and let $D$ be its unique minimum dominating
set. Select two exterior private neighbours of each member of $D$, leaving
two residual vertices $z,w$. Put
\[
 D\cap A=X=\{x_1,\ldots,x_p\},\qquad
 D\cap B=Y=\{y_1,\ldots,y_q\},\qquad p+q=\gamma,
\]
and use $U_i\subseteq B$, $V_j\subseteq A$ for the selected private pairs.
We prove the upper bound
\[
 M_2(\gamma)=\ceil{\frac{\gamma^2}{2}}+5\gamma.
\]

\subsection{Residual vertices in the same part}

Orient the parts so that $z,w\in A$. Each has a nonempty $D$-neighbourhood
contained in $Y$, hence $q\ge1$; $p=0$ is allowed. Call a residual vertex
\emph{singleton-dominated} if it has one neighbour in $D$. Let $S\in\{0,1,2\}$
be the number of such vertices, and let $r_j$ count those whose sole
$D$-neighbour is $y_j$, so that $\sum_jr_j=S$.

For each $i,j$, take the cell with centers $x_i,y_j$, private set $U_i$,
and opposite private set $V_j$ together with the $r_j$ residuals owned by
$y_j$. Every vertex outside this cell remains dominated by
$D\setminus\{x_i,y_j\}$: selected private vertices retain their centers;
singleton residuals outside have a different retained owner; residuals
with at least two $D$-neighbours retain one, because only one $Y$-center
is removed. Lemmas~\ref{lem:six} and~\ref{lem:seven} therefore bound its
optional edges by $2+2r_j$.

These optional edge sets are pairwise disjoint. In addition to the
$2\gamma$ fixed selected-private incidences, count the $S$ singleton owner
edges separately. Each of the remaining $2-S$ residuals has degree at most
$2p+q$: it can meet only $Y$ and the $U$ vertices. There is no edge $zw$.
Consequently
\begin{align}
 \e(G)&\le 2\gamma+S+2pq+2pS+(2-S)(2p+q)\notag\\
       &=2\gamma+2pq+4p+2q-S(q-1)\notag\\
       &\le 2pq+6p+4q.\label{eq:two-same}
\end{align}
This includes $p=0$, with no cells. For $d=p-q$, the final expression is
\begin{equation}\label{eq:two-opt}
 \frac{\gamma^2}{2}+5\gamma+\frac12-\frac{(d-1)^2}{2}.
\end{equation}
If $\gamma$ is even, $d$ is even and $(d-1)^2\ge1$; if $\gamma$ is odd,
$(d-1)^2\ge0$. Thus~\eqref{eq:two-opt} is at most $M_2(\gamma)$.

\subsection{Residual vertices in opposite parts}

Suppose $z\in A$, $w\in B$, and write
\[
 J=N(z)\cap D\subseteq Y,\qquad K=N(w)\cap D\subseteq X.
\]
Both sets are nonempty, so $p,q\ge1$. The edge $zw$ is now possible and
must be counted once.

\paragraph{Both residuals have several dominators.}
Suppose $|J|,|K|\ge2$. Each ordinary six-vertex cell is valid: both
residuals retain a $D$-neighbour after one $X$-center and one $Y$-center
are removed. The cell edges total at most $2pq$. The edges incident with
$z$ other than $zw$ number at most $2p+q$, and those incident with $w$
other than $zw$ number at most $2q+p$. Hence
\begin{equation}\label{eq:two-multi}
 \e(G)\le2\gamma+2pq+(2p+q)+(2q+p)+1=2pq+5\gamma+1.
\end{equation}
For odd $\gamma$, the inequality $2pq\le\floor{\gamma^2/2}$ gives the result.
For even $\gamma$,~\eqref{eq:two-multi} is at most $M_2(\gamma)+1$.
An actual violation would, by integrality, force equality throughout,
$p=q$, and saturation of both residual degree bounds. In particular,
$z$ meets every $Y$-center and every $U$ vertex, while $w$ meets every
$X$-center and every $V$ vertex. For any $i,j$, the set
\[
 (D\setminus\{x_i,y_j\})\cup\{z,w\}
\]
then dominates. The removed centers are covered by $w,z$ respectively,
$U_i$ by $z$, and $V_j$ by $w$; all other private pairs retain their centers.
This second dominating $\gamma$-set excludes the only possible excess.

\paragraph{Exactly one residual has one dominator.}
Interchange sides if necessary so that $J=\{y_{j_0}\}$ and $|K|\ge2$.
For $j\ne j_0$, use ordinary six-vertex cells. For $j=j_0$, include $z$
and use a $(2,3)$ private cell. Outside every such cell, $w$ retains an
$X$-neighbour; any residual not included keeps its owner. Thus the optional
blocks contribute at most $2p(q-1)+4p$. Count $zy_{j_0}$ once, the other
incidences of $w$ by $2q+p$, and $zw$ at most once. This gives
\begin{equation}\label{eq:two-mixed}
 \e(G)\le2\gamma+1+2p(q-1)+4p+(2q+p)+1
       =2pq+5\gamma+2-q.
\end{equation}
If $q\ge2$, this is at most $M_2(\gamma)$. If $q=1$, then $p\ge2$ because
$|K|\ge2$, so $\gamma\ge3$. The difference between $M_2(\gamma)$ and the
right side is $\ceil{\gamma^2/2}-2\gamma+1$, which is zero at $\gamma=3$
and positive for $\gamma\ge4$.

\paragraph{Both residuals have one dominator.}
Write $J=\{y_{j_0}\}$ and $K=\{x_{i_0}\}$. Partition optional edges using
$(2,2)$ cells for $i\ne i_0,j\ne j_0$; $(3,2)$ cells including $w$ for
$i=i_0,j\ne j_0$; $(2,3)$ cells including $z$ for $i\ne i_0,j=j_0$;
and one $(3,3)$ cell containing both residuals for $i=i_0,j=j_0$.
Every residual whose owner is removed is included in the cell, and every
residual outside retains its owner. All selected private vertices outside
also retain their centers. Thus all local uniqueness hypotheses hold.

The edge $zw$ belongs only to the $(3,3)$ cell. The two residual owner
edges are counted once each as fixed incidences, and the optional blocks
are disjoint. The local bounds give
\begin{align}
 \e(G)&\le2\gamma+2+2(p-1)(q-1)+4(q-1)+4(p-1)+6\notag\\
       &=2pq+4\gamma+2\notag\\
       &\le\floor{\gamma^2/2}+4\gamma+2\le M_2(\gamma),\label{eq:two-single}
\end{align}
using $\gamma\ge2$. This covers the small case $p=q=1$ as well.
Together with~\eqref{eq:attain2}, all cases prove the extremal value in
Theorem~\ref{thm:two}.

\subsection{Nonisomorphic equality graphs for even domination number}

Let $\gamma=2k\ge4$. Both $H_2(k,k)$ and $H_2(k+1,k-1)$ attain the
maximum by~\eqref{eq:Hcount}. Their bipartition size pairs are respectively
$\{3k,3k+2\}$ and $\{3k+1,3k+1\}$. A connected bipartite graph has its
bipartition determined up to interchange, so the graphs are nonisomorphic.
The complete equality classification is proved next.

\section{Equality with two residual vertices}\label{sec:two-rigidity}

Assume $\e(G)=M_2(\gamma)$. We retain the private-pair notation of
Section~\ref{sec:two}. No connectedness assumption is used in the classification.

\subsection{Opposite-side residuals above domination number two}

Suppose first that both residuals have several dominators. For odd $\gamma$,
equality in~\eqref{eq:two-multi} forces a balanced split and saturation of
all residual incidences. Replacing any $x_i,y_j$ by $z,w$ gives another
dominating $\gamma$-set, as in the upper-bound proof.

For even $\gamma$, equality requires $p=q=k\ge2$; any other split reduces
$2pq$ by at least two and makes~\eqref{eq:two-multi} strictly smaller than
$M_2(\gamma)$. There is exactly one total deficit in that bound. All
cell deficits and residual-incidence deficits are nonnegative integers.
Thus at most one of the allowable residual edges is missing.
If it is incident with an $X$-center or one of its private vertices,
choose a different $X$-center to remove; if it is incident with a
$Y$-center or its private pair, choose a different $Y$-center. Such
choices exist because $k\ge2$. If the deficit is a cell edge or $zw$,
choose any pair. The set
\[
 (D\setminus\{x_i,y_j\})\cup\{z,w\}
\]
then dominates: the removed centers and their pairs have all needed
residual incidences, and other private pairs retain their centers.
A missing edge to a retained center or its pair is harmless. This again
contradicts uniqueness.

Suppose exactly one residual is singleton-dominated, with
$J=\{y_{j_0}\}$ and $|K|\ge2$. The difference between the target and the
right side of~\eqref{eq:two-mixed} is
\[
 \ceil{\gamma^2/2}-2pq+q-2.
\]
For $q\ge2$ this vanishes only at $p=q=2$: odd $\gamma$ contributes at
least one in the first two terms, and even $\gamma$ requires $p=q$ and
$q=2$. For $q=1$, it vanishes only at $\gamma=3,p=2$.
In either possible split, the two special seven-cells on $x_1,y_{j_0}$
and $x_2,y_{j_0}$ are saturated. By Lemma~\ref{lem:seven-equality}, each
$U_i$ is complete to two members of $W=V_{j_0}\cup\{z\}$.
Choose a common member $c$ of these two two-element subsets of a three-set.
Saturation also makes $w$ adjacent to all $X$-centers, all $V$ vertices,
and $z$. Hence
\[
 \{w,c\}\cup(Y\setminus\{y_{j_0}\})
\]
is a dominating set of size $q+1=\gamma-1$. The vertex $c$ covers both
$U$ pairs and $y_{j_0}$; $w$ covers $X$, all $V$ and $z$; the other
$Y$-centers are selected. This contradicts the domination number.

Finally, if both residuals are singleton-dominated, the gap in
\eqref{eq:two-single} is at least $\gamma-2>0$ for $\gamma\ge3$.
Thus every equality graph with $\gamma\ge3$ has same-side residuals.
These arguments apply to any selected exterior-private pairs, not a
presumed favourable decomposition.

\subsection{Same-side equality with at least two opposite centers}

In~\eqref{eq:two-same}, equality forces $S=0$, complete residual
incidences to every $Y$-center and every $U$ vertex, and two optional
edges in every six-cell. A center edge $x_i y_j$ would allow replacing
$x_i$ by $z$, so all center edges are absent. Each block is a two-edge
star. A star centered at $u\in U_i$ would allow replacing $x_i,y_j$ by
$u,z$; the other residual retains a $Y$-neighbour since $q\ge2$.
Thus every block is a column star.

For a fixed $j$, its chosen column cannot vary with $i$. Otherwise
\[
 \{z\}\cup\{u_i: u_i\in U_i,\ 1\le i\le p\}
        \cup(Y\setminus\{y_j\})
\]
would be another dominating $\gamma$-set, by the argument in
Section~\ref{sec:rigidity}. In this setting $w$ is also covered, since it
retains a $Y$-neighbour. Relabelling the common column of each $V_j$ as
$v_j^0$ gives precisely $H_2(p,q)$.

Equality in~\eqref{eq:two-opt} requires $p-q=1$ for odd $\gamma$ and
$p-q\in\{0,2\}$ for even $\gamma$. These are the splits stated in
Theorem~\ref{thm:two}.

\subsection{Same-side equality with one opposite center}

The allowed splits now give $p=1,2,3$, corresponding to $\gamma=2,3,4$.
Both residuals have the sole center $y$ as owner. Put
$W=V_1\cup\{z,w\}$, of size four. Every $(2,4)$ cell must have six
optional edges.

Its equality pattern has no center edge and consists of $K_{2,3}$ on
$U$ and three columns of $W$. Indeed, if $xy$ is present, uniqueness
excludes a full column, leaving at most four cross-edges and five
optional edges. With $xy$ absent, six cross-edges have row degrees
$(4,2)$ or $(3,3)$. A full row and a neighbour of the other row exclude
$(4,2)$. Different three-element row neighbourhoods give an alternative
pair by selecting one row vertex and the other's neighbour outside that
row. Thus the two rows have the same three neighbours. Conversely,
the unused column and $x$ have disjoint closed neighbourhoods forcing
the center pair, just as in Lemma~\ref{lem:seven-equality}.

Let $c_i$ be the omitted column for $U_i$. If the omissions differ,
the active columns cover $W$, and some column $c$ is active in every
row because $p\le3<4$. Choosing one $u_i\in U_i$, the set
$\{c,u_1,\ldots,u_p\}$ is another dominating $(p+1)$-set: it covers all
centers, all $U$ vertices, and all of $W$. Hence the omitted column is
common, and relabelling the three active columns gives $H_2(p,1)$.

\subsection{The three equality graphs at domination number two}

The remaining possibility has opposite-side residuals and $p=q=1$.
The entire graph is the $(3,3)$ private cell. Write
$U=\{u_0,u_1,u_2\}$ and $V=\{v_0,v_1,v_2\}$, with fixed edges $xU,yV$.
There must be six optional edges.

If $xy$ is absent, there are six cross-edges. With no full row or column,
every row and column would have degree two, and a missing-edge pair
would dominate. By symmetry assume a full row. There is no full column,
or its vertex and the full-row vertex would dominate. The other three
cross-edges occupy the three columns once each. Row degrees $(2,1)$
give a dominating pair, so they must be $(3,0)$. This yields the graph
$F_0$: no edge $xy$, and the six cross-edges $\{u_0,u_1\}\times V$.

If $xy$ is present, the five cross-edges have no full row or column;
the row and column degree lists are $(2,2,1)$. Their graph has maximum
degree two, no isolates, six vertices and five edges. Its only component
types are $P_6$ and $C_4\mathbin{\dot\cup}K_2$. In the path
$u_0v_0u_1v_1u_2v_2$, the pair $\{u_2,v_0\}$ dominates the cell.
Thus only the second type remains. This gives $F_1$: the edge $xy$,
the four cross-edges $\{u_0,u_1\}\times\{v_0,v_1\}$, and $u_2v_2$.

Both graphs have unique minimum dominating pair $\{x,y\}$. In $F_0$,
a dominating pair meets the disjoint sets $N[u_2]=\{u_2,x\}$ and
$N[y]=\{y,v_0,v_1,v_2\}$. Selecting $u_2$ leaves an active $U$ vertex
or a $V$ vertex undominated, forcing $x$; a $V$ vertex instead of $y$
then leaves the other $V$ vertices undominated. In $F_1$, a pair with
one center other than $\{x,y\}$ cannot cover the opposite private set.
A pair inside $U$ misses $y$, and one inside $V$ misses $x$.
An alternative $\{u,v\}$ with $u\in U,v\in V$ would require both
cross-degrees to be two and $uv$ to be a nonedge. All such degree-two
vertices lie in the $C_4$ component, where every opposite pair is
adjacent. A singleton cannot dominate a graph with these part sizes.

Together with the same-side graph $H_2(1,1)$, these are exactly the
three classes in Table~\ref{tab:exceptions}. The degree sequences
separate $F_0,F_1$, and their balanced bipartitions distinguish them
from $H_2(1,1)$. For larger even $\gamma$, the two listed $H_2$ graphs
are nonisomorphic by Section~\ref{sec:two}. All listed graphs are
connected and attain the bound, completing Theorem~\ref{thm:two}.

\begin{table}[htbp]
\centering
\begin{tabular}{lll}
\toprule
Graph & Bipartition sizes & Sorted degree sequence\\
\midrule
$H_2(1,1)$ & $3,5$ & $1,2,3,3,3,4,4,4$\\
$F_0$ & $4,4$ & $1,3,3,3,3,3,4,4$\\
$F_1$ & $4,4$ & $2,2,3,3,3,3,4,4$\\
\bottomrule
\end{tabular}
\caption{The three eight-vertex equality graphs, each with twelve edges.}
\label{tab:exceptions}
\end{table}

\section{Stability in terms of absolute edge deficit}\label{sec:stability}

Let $G\in\mathcal B_{\gamma,1}$ and put
$t=\gamma(\gamma+7)/2-\e(G)$. Theorem~\ref{thm:main} shows $t\ge0$;
if $t=0$, its equality statement gives $d_{\mathrm{edit}}(G,H^*)=0$.
Suppose henceforth $t\ge1$. Selecting two exterior private neighbours
per center leaves $z\in A$, with $p,q,U_i,V_j$ as before.

\subsection{An exact deficit decomposition}

First suppose $z$ has at least two neighbours in $D$. Let $k_{ij}\le2$
be the optional-edge count of the six-cell on $x_i,y_j$. Define
\[
 B=\frac{(p-q)(p-q-1)}2,\qquad
 L=\sum_{i,j}(2-k_{ij}),\qquad R=2p+q-\deg(z).
\]
Each is a nonnegative integer. Complete edge accounting gives the exact
identity
\begin{equation}\label{eq:slack}
 t=B+L+R.
\end{equation}
Here $B$ measures center-split imbalance, $L$ the cell deficits, and $R$
the missing allowable residual incidences.

For an intermediate template, allow $p=0$ and write $T(p,q)$ for the
same adjacency rule as $H(p,q)$ with an empty $X$ collection when
necessary. We still have $q\ge1$. The final balanced template has
both center sides nonempty.

Call row $i$ bad if $z$ misses a member of $U_i$, and column $j$ bad
if $z$ misses $y_j$. Let their counts be $b_X,b_Y$, and let $R_U,R_Y$
count the corresponding missing incidences. Then
\[
 b_X\le R_U,\qquad b_Y\le R_Y,\qquad R_U+R_Y=R.
\]
On a good row, no center edge $x_i y_j$ can occur: replacing $x_i$ by
$z$ would cover its private pair and leave $x_i$ dominated by the
retained $y_j$.

A saturated block in a good row and good column is therefore a
cross-edge star. It cannot be centered at $u\in U_i$, since replacing
$x_i,y_j$ by $u,z$ would dominate. Thus it is a column star.
The selected column is common to all good saturated rows at a fixed
good $j$. Otherwise choose two good rows $i,h$ with opposite selections.
For $u_i\in U_i,u_h\in U_h$, the set
\[
 (D\setminus\{x_i,x_h,y_j\})\cup\{z,u_i,u_h\}
\]
would be another dominating $\gamma$-set. The two selected private
vertices cover their centers and both members of $V_j$, while $z$
covers their private pairs and $y_j$; all other private pairs retain
their centers. This proves coherence even in the presence of bad rows
elsewhere.

Choose that common column for each good $j$, or an arbitrary column
when there is no good saturated row; choose arbitrary columns for bad
$j$. These choices define a labelled $T(p,q)$. Its fixed private edges
already agree with those of $G$. Add the $R$ missing residual incidences.
Every block touching a bad row or column costs at most four edits,
since each graph has at most two optional edges there. The same bound
applies to each deficient good block, of which there are at most $L$.
All remaining blocks agree. Thus
\begin{align}
 d_{\mathrm{edit}}(G,T(p,q))
 &\le R+4L+4(qb_X+pb_Y)\notag\\
 &\le (4\gamma+1)R+4L\notag\\
 &\le (4\gamma+1)t.\label{eq:repair}
\end{align}
The displayed repair uses a common labelling, so allowing arbitrary
bijections only reduces its cost. The argument includes $p=0$.

\subsection{Balancing the template}

Put $a=\ceil{\gamma/2}$, $b=\floor{\gamma/2}$ and $h=|p-a|$.
Writing $p=a+m$ gives $B=2m^2-m$ for even $\gamma$ and $B=2m^2+m$ for
odd $\gamma$. For integer $m$, both are at least $m^2$. Hence
\begin{equation}\label{eq:imbalance-stability}
 h^2\le B\le t.
\end{equation}

Move $h$ center/private triples to the other role, keeping the residual
and all unchanged triples fixed. Keep selected-column labels on
unchanged $Y$-triples. This is a bijection to $T(a,b)=H^*$, and the two
templates agree on all edges whose endpoints both lie outside the
moved triples.

If $p>a$, a moved $X$-triple has degree sum $2q+6$ in the old template
and $2a+5$ in its new $Y$-role. If $p<a$, a moved $Y$-triple has degree
sum $2p+5$, and its new $X$-role has sum $2b+6$. In either case the sum
of old and new degrees is $2\gamma-2h+11$. Counting each changed edge
by an incident moved vertex, with repeated counting allowed for an
upper bound, gives
\begin{equation}\label{eq:balance-repair}
 d_{\mathrm{edit}}(T(p,q),H^*)
 \le h(2\gamma-2h+11)\le(2\gamma+11)\sqrt t.
\end{equation}
By the triangle inequality,~\eqref{eq:repair} and~\eqref{eq:balance-repair}
yield
\[
 d_{\mathrm{edit}}(G,H^*)
 \le(4\gamma+1)t+(2\gamma+11)\sqrt t
 \le(6\gamma+12)t\le12\gamma t,
\]
using the integer condition $t\ge1$ and $\gamma\ge2$.

\subsection{A residual vertex with one owner}

If $z$ has one $D$-neighbour,~\eqref{eq:single} gives
\[
 t\ge B+q-1=\frac{(p-q-1)^2+\gamma-3}{2}
      \ge\floor{\frac{\gamma-2}{2}},
\]
where the last step uses parity. For $\gamma=2,3$ and $t\ge1$, the
trivial distance bound $d_{\mathrm{edit}}(G,H^*)\le\gamma(\gamma+7)$
is at most $12\gamma t$. For $\gamma\ge4$, use
\[
 \gamma+7\le12\floor{\frac{\gamma-2}{2}}
\]
to obtain the same conclusion. This proves the linear-deficit bound
for every graph in the class, without connectedness.
The additional bound
\[
 d_{\mathrm{edit}}(G,H^*)\le\e(G)+\e(H^*)
      =\gamma(\gamma+7)-t
\]
is immediate. Together they establish the upper assertion of
Theorem~\ref{thm:stability}.

\section{Connected obstructions to stronger edit stability}\label{sec:obstruction}

Let $k\ge2$ and $1\le s\le k-1$. Starting from $H(k,k)$, modify each
of the rows $i=1,\ldots,s$ by deleting the $k$ edges $u_i^0v_j^0$,
adding the $k$ edges $x_i y_j$, and deleting $zu_i^0$. Call the result
$G(k,s)$. The same private vertex loses both its selected-column edges
and its residual incidence, becoming a leaf at $x_i$. Its mate $u_i^1$
retains all its edges.

Every change respects the original bipartition. Connectivity is
preserved: $z$ meets all $y_j$, every unmodified $U$ vertex and every
$u_i^1$; each $x_i$ meets $u_i^1$, each $V$ vertex meets its owner, and
the modified leaves meet $x_i$. There are no isolated vertices.
The order is $6k+1$ and the edge count is $2k^2+7k-s$.

\subsection{The unique minimum dominating set}

Every dominating set meets the following pairwise disjoint closed
neighbourhoods:
\[
 \{x_i,u_i^0\}\quad(i\le s),\qquad
 \{x_i,u_i^0,u_i^1\}\quad(i>s),\qquad
 \{y_j,v_j^1\}\quad(1\le j\le k).
\]
They are respectively the closed neighbourhoods of the modified leaves,
the unmodified $X$-centers, and the $V$ leaves. There are $2k$ of them,
so every dominating set has at least $2k$ vertices. The set $D=X\cup Y$
dominates and has that size.

A dominating $2k$-set selects exactly one member of each forced set
and nothing else. Thus $z$, all $v_j^0$, and all modified-row $u_i^1$
are absent. In an unmodified row, selecting a $U$ vertex leaves its mate
undominated. In a modified row, selecting $u_i^0$ instead of $x_i$
leaves $u_i^1$ undominated, since its neighbours are $x_i,z$ and the
$v_j^0$. Therefore every $x_i$ is selected. Selecting $v_j^1$ instead
of $y_j$ then leaves $v_j^0$ undominated. Hence $D$ is the unique minimum
dominating set, and $G(k,s)\in\mathcal B_{2k,1}$ has exact deficit $s$.

\subsection{A degree lower bound for every vertex bijection}

For any bijection between two graphs,
\[
 \sum_v\bigl|\deg_G(v)-\deg_H(\phi^{-1}(v))\bigr|
 \le2|E(G)\mathbin{\triangle}\phi(E(H))|.
\]
\Needspace{7\baselineskip}
The minimum absolute-difference sum between two degree multisets is
achieved by pairing their sorted sequences. Indeed, for $a\le a'$ and
$b\le b'$, uncrossing follows from
\[
 |a-b|+|a'-b'|\le|a-b'|+|a'-b|.
\]
Thus half the sorted degree $L^1$ distance is a lower bound on the
edit distance under arbitrary relabelling.

The degree groups of $H(k,k)$, as (degree, multiplicity), are
\[
 (1,k),\ (2,k),\ (3,k),\ (k+2,2k),\ (2k+1,k),\ (3k,1).
\]
For $G(k,s)$ they are
\[
 (1,k+s),\ (2,k-s),\ (3+s,k),\ (k+2,2k),\
 (2k+1-s,k),\ (3k-s,1).
\]
These groups are nondecreasing for $1\le s\le k-1$; equalities between
groups do not affect the sorted-position calculation. The $L^1$
difference is $s+ks+ks+s=2s(k+1)$. Therefore
\begin{equation}\label{eq:distance-obstruction}
 s(k+1)\le d_{\mathrm{edit}}(G(k,s),H(k,k))\le s(2k+1).
\end{equation}
The upper bound is the explicit modification under the identity
labelling: $s(k+1)$ deletions and $sk$ additions.
This completes Theorem~\ref{thm:stability}.

\subsection{Failure of square-root and dense edit stability}

Take $k=s^2$ with $s\ge2$. The lower bound in
\eqref{eq:distance-obstruction} is $s(s^2+1)$, whereas
$\gamma\sqrt t+t=2s^2\sqrt s+s$. Their ratio tends to infinity.
Hence no absolute constant can bound every graph in the class by
$O(\gamma\sqrt t+t)$ edits.

For a stronger asymptotic consequence, take $s=\floor{k/2}$.
Then $s/(2k^2+7k)\to0$, so the relative edge deficit tends to zero, while
\[
 \liminf_{k\to\infty}
 \frac{d_{\mathrm{edit}}(G(k,\floor{k/2}),H(k,k))}{(6k+1)^2}
 \ge\frac1{72}.
\]
Thus near-optimal edge density does not imply dense edit-distance
closeness, even for connected graphs. In contrast,
Theorem~\ref{thm:stability} guarantees normalized edit closeness when
$t=o(\gamma)$. A single missing residual incidence in the construction
permits changes across an entire row of cells, explaining the linear
rather than square-root dependence on the absolute deficit.

\section{Source comparison and further questions}

At order $3\gamma+1$, Theorem~\ref{thm:main} exceeds the conjectured bound
by $b-1$, as in~\eqref{eq:gap}. At order $3\gamma+2$, the same conjecture
specializes to $2\gamma+2ab+4a+2$. Its summation tail is again zero, both
with the printed cutoff and its construction-consistent correction.
Theorem~\ref{thm:two} exceeds this expression by
\[
 \bigl(2\gamma+2ab+4a+2b\bigr)
       -\bigl(2\gamma+2ab+4a+2\bigr)=2(b-1).
\]
Thus the two sharp boundary families violate the proposed bound for every
$\gamma\ge4$. The individual $13$-vertex and $14$-vertex counterexamples,
and the cutoff issue, were already reported in~\cite{E}. Neither theorem
claims a smallest counterexample over all orders.

For comparison, the published unbalanced-family expression in~\cite{E}
at two extra vertices is $2\gamma+2p'q'+2(2q'+1)$, with $p'+q'=\gamma$.
When $p'=q,q'=p$, its count differs from~\eqref{eq:Hcount} by $2(q-1)$.
At $\gamma=5$, its maximum over positive splits is $36$, whereas
Theorem~\ref{thm:two} gives $38$, attained by $H_2(3,2)$. This compares
the explicit formulas, without asserting priority over every other
possible construction.

Both boundary equality classes are now completely determined. Theorem~\ref{thm:stability}
shows that exact rigidity does not imply dense edit stability. It remains
useful to distinguish other metrics, such as removal of a small vertex set,
or additional assumptions, such as independence of the unique minimum
dominating set. No theorem for those variants is asserted here. Nor does
the present local argument supply a general bound at orders $3\gamma+r$
with $r\ge3$: the residual vertices' simultaneous domination obligations
must still be retained.

\appendix
\section{Supplementary exact checks}\label{app:checks}

Independent exact checkers enumerate all $1696$ optional-edge patterns
on $xy$ and $U\times V$, testing every vertex subset of size at most two.
Table~\ref{tab:cells} records their results. These computations supplement
the analytic proofs.

\begin{table}[htbp]
\centering
\begin{tabular}{rrrr}
\toprule
$(a,b)$ & Optional patterns & Admissible patterns & Maximum optional edges\\
\midrule
$(2,2)$ & 32 & 14 & 2\\
$(2,3)$ & 128 & 58 & 4\\
$(2,4)$ & 512 & 254 & 6\\
$(3,3)$ & 1024 & 548 & 6\\
\bottomrule
\end{tabular}
\caption{Admissibility means that $\{x,y\}$ is the only dominating set of size at most two.}
\label{tab:cells}
\end{table}

The first-boundary low-$\gamma$ equality reduction was checked on all
$128$ seven-vertex and $16\,384$ ten-vertex skeletons, including
subextremal patterns. Each has three labelled equality patterns.
For the second boundary, all $15$ maximum $(3,3)$ patterns form the
two balanced eight-vertex classes, of labelled multiplicities $6$ and
$9$, under independent full-graph isomorphism checks. Among the $164$
reduced same-side patterns through $\gamma=5$, exactly $20$ survive,
four at each tested split; those survivors map to the stated $H_2$
representatives. The others have alternative dominating sets.
All $9$ and $576$ mixed-owner candidates at $\gamma=3,4$ have the
smaller dominating set used in the proof. These are reduced-skeleton
checks, not unrestricted graph censuses or counts of isomorphism classes.

All subsets were tested for the ten $H_2(p,q)$ graphs with
$p,q\ge1$ and $2\le p+q\le5$; $H_1$ was checked through $\gamma=6$.
Explicit isomorphisms recover the prior $13$- and $14$-vertex certificates
in~\cite{E}. Stability checks independently verify domination, forced
neighbourhoods, sorted degrees and $819$ balancing bijections through
$\gamma=40$; reduced-graph checks also confirm the repair step.
The accompanying package supplies the checkers, results, source hashes
and independent audits. No Lean formalization is claimed.

\begin{thebibliography}{9}
\bibitem{FRV}
M.~Fischermann, D.~Rautenbach and L.~Volkmann,
\emph{Maximum graphs with a unique minimum dominating set},
Discrete Mathematics \textbf{260} (2003), 197--203.

\bibitem{FS}
M.~Fraboni and N.~Shank,
\emph{Maximum graphs with unique minimum dominating set of size two},
Australasian Journal of Combinatorics \textbf{46} (2010), 91--99.
\url{https://ajc.maths.uq.edu.au/pdf/46/ajc_v46_p091.pdf}.

\bibitem{E}
J.~Erlbacher,
\emph{Demonstrandum---wave-1 verified artifacts: Koch--Narayan counterexample package},
public research artifacts, 2026.
\href{https://github.com/demonstrandum-research/artifacts/blob/94db9ed50d48a57aae5ccb72e6a95a2b8f8f39d3/problems/p2-factory/kills/koch-narayan/WRITEUP.md}{Pinned write-up, commit \texttt{94db9ed50d48} (13 June 2026)}.

\bibitem{KN}
G.~Koch and D.~Narayan,
\emph{Maximal bipartite graphs with a unique minimum dominating set},
arXiv:2511.01719v1, 3 November 2025.
\url{https://arxiv.org/abs/2511.01719v1}.
\end{thebibliography}
\end{document}
